\documentclass[runningheads]{llncs}

% Springer CCIS/LNCS proceedings format for SOICT 2026.
% Do not add geometry, fancyhdr, titlesec, setspace, or custom font packages:
% the llncs class controls the required page size, margins, headings, and fonts.
\usepackage[T1]{fontenc}
\usepackage{amsmath,amssymb}
\usepackage{graphicx}
\usepackage{booktabs}
\usepackage{url}
\usepackage[hidelinks]{hyperref}

\begin{document}

\title{LegalSCM-RAG: Causal Path Retrieval and Selective Structural Verification for Vietnamese Multi-Hop Legal Question Answering}
\titlerunning{LegalSCM-RAG}

% SOICT 2026 uses single-blind review, so author names and affiliations
% must remain in the submission manuscript.
\author{Nguyen Thi Nhung, Nguyen Quoc Tuan, Nguyen Thi Kim Anh}
% \authorrunning{N. T. Nhung}

\institute{Hanoi University of Science and Technology, Hanoi, Vietnam\\
\email{ }}

\maketitle
% SOICT 2026 requires submission PDFs without page numbers.
\pagestyle{empty}

\begin{abstract}
Large language models and Retrieval-Augmented Generation (RAG) remain limited on legal questions that require linking multiple statutory condition--consequence relations. We propose \textbf{LegalSCM-RAG}, which combines causal-path retrieval with a deterministic three-valued normative rule model and a selective commit gate. The formulation distinguishes bounded node-deletion reachability, used only as a topological diagnostic, from the hard intervention $do(X_M=0)$, which disables mechanisms producing mediator $M$ and recomputes the rule fixed point. Conjunction is represented within a rule and alternative disjunctive support by multiple rules with the same consequent. On 250 graph-derived Vietnamese Penal Code questions, Full LegalSCM attains 96.80\% decision accuracy, compared with 90.80\% for Causal Path RAG and 55.60\% for Vanilla RAG. The paired improvement over Causal Path RAG is 6.0 percentage points (95\% CI: 2.8--9.6; Holm-adjusted $p=0.0029$), while Answer F1, ROUGE-L, and Citation F1 do not differ significantly. A route audit shows 230 factual-validation cases and 20 direct-edge checks, with no executed mediator intervention in this batch. The experiment therefore supports the decision-level benefit of selective structural verification, but does not yet establish the effectiveness of $do$-based mediator reasoning.

\keywords{Retrieval-Augmented Generation \and Multi-hop Question Answering \and Three-Valued Logic \and Structural Verification \and Legal Question Answering \and Vietnamese Legal NLP}
\end{abstract}

\section{Introduction}
\label{sec:introduction}

Large Language Models (LLMs) have demonstrated strong performance across a wide
range of question-answering tasks. In the legal domain, however, a reliable
answer requires more than fluent generation: the system must identify relevant
provisions and reconstruct the reasoning chain connecting conditions,
intermediate events, and legal consequences
\cite{chalkidis2020legalbert,guha2023legalbench}.

RAG augments an LLM with evidence retrieved before generation
\cite{lewis2020retrieval,gao2024ragsurvey}. Conventional RAG is nevertheless
primarily driven by semantic similarity and often treats passages as
independent evidence. This is insufficient for multi-hop legal questions whose
conclusion depends on a sequence of connected condition--consequence rules.

We propose \textbf{LegalSCM-RAG}, which combines directed legal-graph retrieval
with question-aware structural verification. The verifier routes direct claims
to grounded edge checking, ordinary causal questions to three-valued factual
rule inference, and explicit mediator-removal questions to hard intervention.
A conservative commit gate uses a verifier decision only when grounding, path
coverage, confidence, and state-consistency requirements are satisfied;
otherwise it preserves the Causal Path RAG context and LLM decision.

The main contributions are:
\begin{itemize}
    \item a three-valued fixed-point formulation with an explicit sparse
    body/head incidence encoding, in which conjunction occurs within a rule and
    alternative disjunctive support across rules;
    \item a formal and implementation-level separation between graph
    node-deletion reachability and mechanism-replacement $do(X_M=0)$; and
    \item a controlled four-system evaluation with paired significance tests,
    abstention analysis, and route/commit auditing.
\end{itemize}

On the current benchmark, selective verification significantly improves
structured decision accuracy over Causal Path RAG, but not free-form answer or
citation quality. Importantly, the evaluated batch contains factual-validation
and direct-edge routes but no executed mediator intervention. We therefore
report the measured gain as evidence for the selective verifier as a whole,
not as an empirical validation of structural counterfactual intervention.


% =========================================================
% II. RELATED WORK
% =========================================================
\section{Related Work}
\label{sec:related_work}

\subsection{Retrieval-Augmented Generation}

Retrieval-Augmented Generation (RAG) combines a language model with an external
retrieval component to provide relevant evidence prior to answer generation
\cite{lewis2020retrieval,gao2024ragsurvey}. Dense retrieval methods such as DPR
\cite{karpukhin2020dense} improve semantic document retrieval by representing
queries and passages in a shared embedding space. Subsequent work has explored
the aggregation of multiple retrieved passages, adaptive retrieval and
self-reflection, and hierarchical retrieval structures
\cite{izacard2021fid,asai2024selfrag,sarthi2024raptor}. However, most
conventional RAG systems still provide retrieved evidence primarily as textual
context rather than as an explicit reasoning structure.

This paradigm is effective for questions that can be answered using one or a
small number of directly relevant passages. Its effectiveness is more limited,
however, for tasks that require information to be connected across multiple
reasoning steps.

\subsection{Multi-hop Question Answering}

Multi-hop Question Answering requires a system to integrate information from
multiple pieces of evidence in order to derive a final answer. Benchmarks such
as HotpotQA \cite{yang2018hotpotqa} and MuSiQue \cite{trivedi2022musique}
demonstrate that retrieving individually relevant documents is not sufficient;
a system must also identify the reasoning chain that connects the supporting
evidence.

This requirement is particularly important in the legal domain, where a legal
conclusion may depend on multiple conditions, exceptions, and intermediate
consequences. Explicitly representing such reasoning structures may therefore
help reduce unsupported reasoning transitions and improve the traceability of
generated conclusions.

\subsection{Graph-based and Causal RAG}

Recent research has explored graph-based representations as a means of
organizing knowledge and supporting multi-hop retrieval. GraphRAG leverages
graph structures to expand retrieval and aggregate information across related
entities and passages \cite{edge2024graphrag}. Think-on-Graph further
demonstrates that explicit traversal over a knowledge graph can support
multi-step reasoning and provide traceable reasoning paths
\cite{sun2024thinkongraph}.

More closely related to our work, CausalRAG incorporates causal relations into
retrieval and reasoning rather than relying solely on semantic similarity
\cite{wang2025causalrag}. CC-RAG similarly constructs causal chains to support
structured multi-hop inference over domain-specific corpora
\cite{parekh2025ccrag}. This line of work suggests that explicit relations
among events can provide useful structural signals for multi-hop reasoning.

In the legal setting considered in this study, however, these relations are
interpreted as directed condition--consequence relations derived from statutory
text. They should therefore be understood as a structured representation of
legal dependencies rather than as causal effects estimated from observational
data.

\subsection{Counterfactual Reasoning and Legal QA}

Counterfactual reasoning provides a mechanism for examining whether a
conclusion would remain valid if an intermediate condition or event were
modified or removed. In multi-hop question answering, prior work on
counterfactual questions has shown that a model may arrive at the correct
answer while relying on an inappropriate, incomplete, or disconnected
reasoning chain \cite{guo2023counterfactual}.

In the legal domain, benchmarks such as LexGLUE, CaseHOLD, and LegalBench, as
well as domain-specific models such as LEGAL-BERT, have highlighted both the
distinctive characteristics of legal language and the continuing difficulty of
structured legal reasoning
\cite{chalkidis2022lexglue,zheng2021casehold,guha2023legalbench,chalkidis2020legalbert}.
These limitations motivate approaches that combine evidence retrieval with more
explicit representations of legal relations and reasoning structure.

\subsection{Three-Valued and Linear-Algebraic Logic Programming}

Three-valued logic distinguishes true, false, and an undefined or unknown state
\cite{kleene1952metamathematics}. Fitting's Kripke--Kleene semantics provides a
fixed-point account for logic programs under partial information
\cite{fitting1985kripke}. More recent work encodes two- and three-valued logic
program semantics in vector spaces \cite{sato2020threevalued}, while matrix and
third-order tensor representations support definite, disjunctive, and normal
logic programs \cite{sakama2021tensor}. These studies motivate our algebraic
formulation. Our runtime remains a symbolic sparse-rule engine: it uses AND
within one rule and OR across independently activated rules with the same
consequent rather than executing tensor multiplication.

\subsection{Research Gap}

Conventional RAG primarily focuses on retrieving semantically relevant
evidence, whereas graph-based approaches emphasize structural connections
among entities, events, or passages. Nevertheless, relatively limited research
has investigated legal QA frameworks that jointly incorporate the following
capabilities:

\begin{itemize}
    \item representing legal rules as directed condition--consequence relations;
    \item retrieving and ranking multi-hop reasoning paths;
    \item assessing the necessity of intermediate events through counterfactual
    analysis; and
    \item using verification outcomes to refine the evidence provided to the
    answer-generation model.
\end{itemize}

This work addresses this gap by integrating causal graph retrieval, multi-hop
path reasoning, and counterfactual verification within a unified RAG framework
for Vietnamese multi-hop legal question answering.

% =========================================================
% III. PROBLEM FORMULATION
% =========================================================
\section{Problem Formulation}
\label{sec:problem_formulation}

We consider multi-hop legal question answering, where a conclusion may require
connecting several statutory conditions, intermediate events, and legal
consequences. Let the input question be $q\in\mathcal{Q}$. Legal knowledge is
represented by a directed graph $G=(V,E)$ and an associated rule set
$\mathcal{R}$. Each node $v_i\in V$ denotes a normalized legal event, while an
edge summarizes one or more statutory condition--consequence rules. The rule
set, rather than graph reachability alone, defines the structural inference
semantics.

Given $q$, the system retrieves relevant events and constructs candidate
reasoning paths
\begin{equation}
\mathcal{P}(q)=\{P_1,P_2,\ldots,P_K\},
\end{equation}
where
\begin{equation}
P_i=v_0\xrightarrow{r_1}v_1\xrightarrow{r_2}\cdots
\xrightarrow{r_h}v_h.
\end{equation}
The primary path is selected by
\begin{equation}
P^*=\arg\max_{P_i\in\mathcal{P}(q)}S(P_i,q),
\end{equation}
where $S(P_i,q)$ combines semantic and structural relevance.

\subsection{Three-Valued Rule Semantics}

Each event is associated with a structural variable
\begin{equation}
X_i\in\{1,0,\bot\},
\end{equation}
where $1$, $0$, and $\bot$ denote \textsc{true}, \textsc{false}, and
\textsc{unknown}. The third value represents evidence that is not established
or is conflicting; it is not interpreted as a numeric value between zero and
one. This design follows three-valued fixed-point semantics for logic programs
\cite{kleene1952metamathematics,fitting1985kripke} and their linear-algebraic
extensions \cite{sato2020threevalued}.

To encode the state without assigning a numeric value to $\bot$, we use two
support channels
\begin{equation}
t_i=\mathbb{I}[X_i=1],\qquad
f_i=\mathbb{I}[X_i=0].
\label{eq:two_channel_state}
\end{equation}
Thus, $X_i=\bot$ gives $(t_i,f_i)=(0,0)$. Conflicting positive and negative
support is conservatively projected to $\bot$ by the runtime inference engine.

Let $\mathbf{B}^{+},\mathbf{B}^{-}\in\{0,1\}^{|\mathcal{R}|\times |V|}$ be
body-incidence matrices for literals requiring true and false states. For rule
$r$, let $k_r=\sum_i(B^{+}_{ri}+B^{-}_{ri})$. Its conjunctive body is satisfied
iff every literal has the required settled state:
\begin{equation}
b_r(\mathbf{t},\mathbf{f})=
\mathbb{I}\!\left[
\sum_i B^{+}_{ri}t_i+
\sum_i B^{-}_{ri}f_i=k_r
\right].
\label{eq:conjunctive_body}
\end{equation}
If $e_r$ indicates that at least one encoded exception is satisfied, the rule
activation is $a_r=b_r(1-e_r)$. Let $\mathbf{H}^{+}$ and $\mathbf{H}^{-}$ be
head-incidence matrices. Positive support for event $i$ is aggregated as
\begin{equation}
d_i^{+}=\mathbb{I}\!\left[\sum_r H^{+}_{ir}a_r\geq 1\right],
\label{eq:disjunctive_support}
\end{equation}
with an analogous expression for negative support. Consequently, literals
inside one rule are conjunctive, whereas independently activated rules with the
same head are alternative, disjunctive mechanisms. For example,
$(A\land B)\rightarrow Y$ is one rule with two body entries, while
$(A\lor B)\rightarrow Y$ is normalized into $A\rightarrow Y$ and
$B\rightarrow Y$. Disjunctive heads are outside the current rule schema.
This body/head incidence view is the sparse matrix counterpart of
linear-algebraic and tensor-space representations of definite, disjunctive,
and normal logic programs \cite{sakama2021tensor}.

The factual state is the fixed point
\begin{equation}
\mathbf{x}^{F}
=
\operatorname{Fix}(\Phi_{\mathcal{R}},\mathbf{x}^{(0)}),
\label{eq:factual_fixed_point}
\end{equation}
where $\Phi_{\mathcal{R}}$ activates the encoded rule mechanisms. The matrices
above define a formal algebraic representation. The implementation evaluates
the corresponding sparse rules symbolically and iterates to a fixed point; it
does not perform dense matrix multiplication. The engine supports multiple
body literals, but the legacy rule collection used in the reported experiment
maps each record to one positive antecedent and one positive consequent.
Therefore, conjunction is an engine-level capability and is not empirically
validated by the present benchmark.

\subsection{Topological Ablation and Structural Intervention}

Let $M(P^*)=\{v_1,\ldots,v_{h-1}\}$ denote the intermediate events of the
primary path. We distinguish two non-equivalent operations. First, a
node-deletion diagnostic forms the induced subgraph
\begin{equation}
G^{-M}=G[V\setminus\{M\}]
\end{equation}
and checks bounded reachability
\begin{equation}
\rho_{\mathrm{top}}(s,Y\mid M)
=
\mathbb{I}[Y\text{ is reachable from }s\text{ in }G^{-M}].
\label{eq:topological_ablation}
\end{equation}
This operation reports whether an alternative graph path is found within the
search limits. It is used only as a topology diagnostic or fallback and is not
the LegalSCM counterfactual definition.

Second, for an explicit mediator-removal query, LegalSCM applies the hard
intervention
\begin{equation}
do(X_M=0),
\label{eq:problem_do}
\end{equation}
where $0\equiv\mathrm{FALSE}$. The variable and the remaining rule structure
are preserved; every rule whose head is $X_M$ is disabled and the mechanism for
$X_M$ is replaced by the constant assignment $X_M:=0$. Equivalently,
\begin{equation}
\mathcal{R}^{do(X_M=0)}
=
\{r\in\mathcal{R}\mid \operatorname{head}(r)\neq X_M\}.
\end{equation}
The counterfactual world is recomputed as
\begin{equation}
\mathbf{x}^{CF}_{M}
=
\operatorname{Fix}\!\left(
\Phi_{\mathcal{R}^{do(X_M=0)}},
\mathbf{x}^{(0)}\cup\{X_M=0\}
\right).
\label{eq:cf_fixed_point}
\end{equation}
For a target $Y$, $M$ is necessary relative to the encoded normative model when
$x_Y^{F}=1$ and $x_{Y,M}^{CF}=0$; it is non-necessary when
$x_{Y,M}^{CF}=1$, and remains indeterminate when either world is unknown.
This is a hard intervention over the encoded rule model, without an exogenous
abduction layer, and must not be equated with deleting $M$ from $G$.

Finally, the system predicts a query-relative decision $z\in\mathcal{Z}$,
where
\begin{equation}
\begin{aligned}
\mathcal{Z}=\{&\textsc{supported},\textsc{uncertain},\\
&\textsc{reject-direct-claim}\},
\end{aligned}
\end{equation}
and generates an answer $a$ with statutory citations $C$:
\begin{equation}
f:(q,G,\mathcal{R})\rightarrow(P^*,z,\mathcal{E}^*,a,C),
\end{equation}
where $\mathcal{E}^*$ is the refined evidence supplied to the generation model.
The decision labels are distinct from both the event-state domain
$\{1,0,\bot\}$ and internal mediator statuses such as necessary or
non-necessary.

% =========================================================
% IV. LEGALSCM-RAG
% =========================================================
\section{LegalSCM-RAG}
\label{sec:legalscm_rag}

We propose \textit{LegalSCM-RAG}, a retrieval-augmented generation framework
for multi-hop legal question answering that integrates retrieval over a
directed condition--consequence graph, multi-hop path reasoning, and selective
verification through a Legal Structural Causal Model (\textit{LegalSCM}).

Given an input question $q$, the overall pipeline is
\begin{equation}
\begin{aligned}
q
&\rightarrow
\text{Causal Retrieval}
\rightarrow
\mathcal{P}(q)
\rightarrow
P^*\\
&\rightarrow
\text{Question-Aware Verification}
\rightarrow
\mathcal{E}^*
\rightarrow
(a,C),
\end{aligned}
\label{eq:overall_pipeline}
\end{equation}
where $P^*$ is the primary reasoning path, $\mathcal{E}^*$ is the selected
evidence, $a$ is the answer, and $C$ is the set of cited provisions.

\subsection{Legal Graph-Based Retrieval}

The system embeds the question and retrieves semantically relevant legal events
and rules from causal memory. Retrieved events serve as seeds for graph
expansion. Starting from each seed, the system traverses directed
condition--consequence relations to construct
\begin{equation}
\mathcal{P}(q)=\{P_1,P_2,\ldots,P_K\},
\end{equation}
where
\begin{equation}
P_i=v_0\xrightarrow{r_1}v_1\xrightarrow{r_2}\cdots
\xrightarrow{r_h}v_h.
\end{equation}
Unlike conventional semantic RAG, the retrieval unit includes legal events and
the directed rule relations connecting them, rather than only independent text
passages.

\subsection{Reasoning Path Ranking}

Candidate paths are ranked using semantic and structural signals:
\begin{equation}
S(P_i,q)=
\alpha S_{\mathrm{sem}}+
\beta S_{\mathrm{graph}}+
\gamma S_{\mathrm{seed}}+
\delta S_{\mathrm{path}}.
\end{equation}
The primary path is
\begin{equation}
P^*=\arg\max_{P_i\in\mathcal{P}(q)}S(P_i,q).
\end{equation}
Rules and provisions associated with $P^*$ are passed to verification as
structured evidence.

\subsection{LegalSCM Fixed-Point Inference}

LegalSCM is a deterministic three-valued normative rule model. Each event is in
$\{\mathrm{TRUE},\mathrm{FALSE},\mathrm{UNKNOWN}\}$. A rule has the form
\begin{equation}
X_1\land\cdots\land X_k\rightarrow Y.
\end{equation}
It fires only when all condition literals have their required settled states
and no encoded exception is active. Conditions within one rule are therefore
conjunctive. Multiple independently activated rules deriving the same state of
$Y$ are alternative, disjunctive support. Opposite supports are retained as a
conflict and mapped to an effective \textsc{unknown} state. The rule operator
is iterated until a fixed point is reached, as formalized in
Section~\ref{sec:problem_formulation}.

The canonical engine supports multiple conditions and exceptions, but the
legacy records used in the reported experiment instantiate one positive
condition and one positive effect per rule. Consequently, the present results
do not constitute an empirical evaluation of multi-condition conjunction or
explicit exception handling.

LegalSCM models normative dependencies extracted from statutes. It neither
discovers causal relations from observational data nor estimates statistical
causal effects.

\subsection{Hard Intervention and Topological Diagnostic}

For an explicit query about removing an intermediate event $M$, LegalSCM
performs a hard intervention relative to the encoded rule model
\cite{pearl2009causality}:
\begin{equation}
do(X_M=0),\qquad 0\equiv\mathrm{FALSE}.
\end{equation}
The intervention fixes $X_M$ to false, blocks every rule mechanism whose effect
is $X_M$, and recomputes the fixed point. If the factual target $Y$ is true and
\begin{equation}
Y_{do(X_M=0)}=\mathrm{FALSE},
\end{equation}
then $M$ is necessary under the encoded context. If
$Y_{do(X_M=0)}=\mathrm{TRUE}$, an alternative mechanism supports the target.
An unknown state in either world yields an indeterminate result.

This operation is distinct from the auxiliary topological ablation that forms
$G^{-M}=G[V\setminus\{M\}]$ and searches for a bounded alternative path. The
latter removes the node and its incident edges and is retained only as a
diagnostic or fallback. It cannot be interpreted as a Pearl-style intervention.
Moreover, LegalSCM performs no exogenous-state abduction; the result is
model-relative necessity in the supplied normative context, not unit-level
actual causation.

\subsection{Question-Aware Verification}

The verifier selects one of three routes from the query. A direct claim
$A\rightarrow C$ is checked against conservatively grounded graph endpoints.
A standard legal question receives factual LegalSCM validation of the primary
path without an intervention. Only an explicit mediator-removal or necessity
query requests $do(X_M=0)$ and compares factual and intervened worlds. This
routing prevents a topological direct-edge test or ordinary factual validation
from being mislabeled as a structural intervention.

\subsection{Selective Verification}

The system distinguishes event states
$\{\mathrm{TRUE},\mathrm{FALSE},\mathrm{UNKNOWN}\}$, internal intervention
statuses such as necessary or non-necessary, and final query labels
\begin{equation}
\mathcal{Z}=\{\textsc{supported},\textsc{reject-direct-claim},
\textsc{uncertain}\}.
\end{equation}
A verifier label becomes authoritative only when confidence, endpoint or
mediator grounding, primary-path coverage, and state-consistency checks pass:
\begin{equation}
z=
\begin{cases}
z_{\mathrm{SCM}}, & G_{\mathrm{commit}}=1,\\[3pt]
z_{\mathrm{LLM}}, & G_{\mathrm{commit}}=0.
\end{cases}
\end{equation}
When the gate fails, the verifier does not force an \textsc{uncertain} label.
Instead, the unchanged Causal Path RAG evidence is supplied to the common
generation model. Full LegalSCM is therefore a selective hybrid rather than an
SCM-authoritative classifier for every sample.

\subsection{Evidence Refinement and Answer Generation}

After verification, the system constructs
\begin{equation}
\mathcal{E}^*=\{r_j\mid r_j\text{ supports }P^*\}
\end{equation}
and provides
\begin{equation}
I_q=(q,P^*,z,\mathcal{E}^*)
\end{equation}
to the language model, which generates
\begin{equation}
(a,C)=f_{\theta}(I_q).
\end{equation}
This separation makes retrieval, rule inference, the selected verification
route, and generation independently auditable.

% =========================================================
% V. EXPERIMENTAL SETUP
% =========================================================
\section{Experimental Setup}
\label{sec:experimental_setup}

\subsection{Evaluation Data}

Experiments use 250 multi-hop questions constructed in this study from the
Vietnamese Penal Code and the derived legal graph. Each sample contains a
question, relevant rules and events, a gold reasoning path, a reference
decision, a reference answer, and legal citations. Of the 250 samples, 236 have
a linear gold path and are used for path evaluation.

The decision labels are imbalanced: 230 samples are \textsc{supported} and 20
are \textsc{reject-direct-claim}; no gold sample is \textsc{uncertain}. The 20
negative samples ask whether an endpoint causes another endpoint directly,
without the intermediate event of a supported two-hop chain. They are therefore
\emph{direct-claim counterexamples}. They do not explicitly request
$do(X_M=0)$ and are routed to grounded direct-edge topology. The reported batch
contains no explicit mediator-intervention query; this distinction is retained
in the result interpretation.

Because both the questions and gold paths are derived from the same structured
legal resource used to build the retrieval graph, the experiment measures
controlled internal validity rather than external expert-level legal validity.

\subsection{Compared Systems}

We evaluate four configurations:
\begin{itemize}
    \item \textbf{Closed-book LLM}: Qwen3-8B receives only the question under
    the shared verification-oriented output contract. The prompt requires the
    model to return \textsc{uncertain} instead of fabricating unsupported
    evidence. This is a conservative abstaining baseline, not an unrestricted
    estimate of general legal ability.
    \item \textbf{Vanilla RAG}: retrieves statutory provisions by semantic
    similarity before generation.
    \item \textbf{Causal Path RAG}: retrieves rules, events, and multi-hop paths
    without LegalSCM verification.
    \item \textbf{Full LegalSCM}: adds question-aware LegalSCM verification and
    the selective commit gate to the same Causal Path RAG pipeline.
\end{itemize}

All systems use the same model, JSON output schema, system-prompt rules, context
limit, and decoding parameters. The prompt policy differs only in whether no
context, retrieved evidence, or a committed authoritative verifier decision is
available.

\subsection{Experimental Configuration}

Qwen3-8B is executed locally through Ollama. Table~\ref{tab:experimental_config}
shows the common settings.
\begin{table}[t]
\centering
\caption{Main experimental configuration.}
\label{tab:experimental_config}
\small
\begin{tabular}{lc}
\hline
\textbf{Parameter} & \textbf{Value} \\
\hline
LLM & Qwen3-8B \\
Embedding model & BAAI/bge-m3 \\
Temperature & 0.1 \\
Random seed & 42 \\
Maximum output tokens & 1200 \\
Context window & 32768 \\
Maximum context characters & 18000 \\
Retrieval $K$ & 5 \\
Maximum graph hops & 2 \\
\hline
\end{tabular}
\end{table}
BGE-M3 \cite{chen2024bgem3} encodes questions and legal knowledge, and FAISS
performs nearest-neighbor search over causal memory.

\subsection{Metrics and Run Audit}

Retrieval is evaluated with Recall@5, MRR, and MAP for each applicable
component. Top-1 Exact Path and Edge F1 evaluate the selected linear path;
Oracle Exact Path diagnoses whether a gold path occurs anywhere in the
candidate set. Oracle performance is not a deployable score.

Decision Accuracy is exact agreement among
\textsc{supported}, \textsc{reject-direct-claim}, and \textsc{uncertain}. We
also report Macro-F1 and Balanced Accuracy because the gold labels are
imbalanced. Token-F1 and ROUGE-L evaluate answers, and precision, recall, and F1
evaluate citations. Missing or failed predictions receive zero on shared
end-to-end metrics.

For every run we audit prediction coverage, failures, generation signatures,
query route, commit-gate status, fallback status, and whether a structural
intervention was actually executed. This prevents a direct-edge or factual
route from being counted as a $do$-based counterfactual result.

\subsection{Statistical Significance Testing}

All comparisons are paired on the same 250 questions. We estimate 95\%
confidence intervals using 5,000 bootstrap samples. Decision Accuracy uses the
exact McNemar test; Answer F1, ROUGE-L, and Citation F1 use 5,000 paired
randomization sign-flips. Holm correction controls multiple comparisons, with
random seed 42.

% =========================================================
% VI. RESULTS AND ANALYSIS
% =========================================================
\section{Results and Analysis}
\label{sec:results_analysis}

\subsection{Overall Results and the Closed-Book Baseline}

All four runs produced 250/250 successful predictions with no pipeline errors
and shared the same model and generation settings. Table~\ref{tab:main_results}
reports end-to-end quality; retrieval/path cells are discussed only for systems
that expose those components.

\begin{table}[t]
\centering
\caption{Overall performance on 250 questions (percent except time).}
\label{tab:main_results}
\scriptsize
\setlength{\tabcolsep}{1.35pt}
\begin{tabular}{@{}lrrrrrrr@{}}
\hline
\textbf{Method}
& \textbf{Acc.}
& \textbf{B.Acc.}
& \textbf{M-F1}
& \textbf{Ans.F1}
& \textbf{R-L}
& \textbf{Cit.F1}
& \textbf{Sec.}\\
\hline
Closed-book LLM & 0.00 & 0.00 & 0.00 & 12.98 & 11.27 & 0.00 & 0.841\\
Vanilla RAG & 55.60 & 59.89 & 67.60 & 23.28 & 21.45 & 49.37 & 2.308\\
Causal Path RAG & 90.80 & 88.15 & 84.36 & \textbf{36.53} & \textbf{35.82} & 58.70 & 2.107\\
Full LegalSCM & \textbf{96.80} & \textbf{95.98} & \textbf{93.31} & 36.32 & 35.73 & \textbf{59.30} & 3.324\\
\hline
\end{tabular}
\end{table}

The 0.00\% closed-book decision score is not a failed run or an evaluator
parsing error. All 250 successful generations returned the valid label
\textsc{uncertain}. Since the gold set contains 230 \textsc{supported}, 20
\textsc{reject-direct-claim}, and no \textsc{uncertain} item, exact-label
accuracy is zero. The result therefore measures \emph{complete abstention under
this conservative shared prompt/model configuration}; it must not be read as
zero general legal competence. Conversely, an always-\textsc{supported}
classifier would obtain 92\% raw accuracy on this imbalanced set, which is why
Balanced Accuracy and Macro-F1 are reported alongside Accuracy.

\subsection{Paired Effect of Selective Verification}

Full LegalSCM answers 242 questions correctly versus 227 for Causal Path RAG.
Among paired outcomes, Full alone is correct on 17 samples, Causal Path RAG
alone on 2, both are correct on 225, and both are wrong on 6. The accuracy gain
is therefore 6.0 percentage points, with paired bootstrap 95\% CI
$[2.8,9.6]$ points, exact McNemar $p=0.000729$, and Holm-adjusted
$p=0.002914$. Macro-F1 increases descriptively by 8.94 points, from 84.36\% to
93.31\%.

The gain is decision-specific. Relative to Causal Path RAG, Full changes Answer
F1 by $-0.21$ points (95\% CI $[-1.83,1.43]$), ROUGE-L by $-0.09$ points
(95\% CI $[-1.71,1.56]$), and Citation F1 by $+0.59$ points (95\% CI
$[-0.72,2.03]$). All three Holm-adjusted $p$-values equal 1.0. Thus the data
support improved structured decisions, not a general improvement in answer
wording or citations.

\subsection{Route, Commit, and Error Audit}

Table~\ref{tab:route_audit} reports independent runtime counters. The query
router sends 230 samples to factual LegalSCM validation and 20 direct claims to
grounded edge checking. No query requests the mediator-intervention route.

\begin{table}[t]
\centering
\caption{Full LegalSCM route and commit audit.}
\label{tab:route_audit}
\small
\begin{tabular}{lr}
\hline
\textbf{Audit quantity} & \textbf{Count}\\
\hline
Factual LegalSCM route & 230\\
Grounded direct-edge route & 20\\
Verifier committed & 230\\
No-primary-path safe fallback & 20\\
Executed $do(X_M=0)$ intervention & 0\\
\hline
\end{tabular}
\end{table}

All 17 Full-only gains occur when the verifier commits: 14 arise from factual
rule validation and 3 from direct-edge topology. In the 20 safe-fallback cases,
Full and Causal Path RAG make the same decision on 19 samples; 13 are jointly
correct and 6 jointly wrong. The remaining case is correct only for Causal Path
RAG. Hence the fallback contributes no Full-only gain.

Full makes eight decision errors. Seven occur after no-primary-path fallback,
primarily because first-hop retrieval does not expose a usable primary path.
The remaining error is a direct-claim endpoint-grounding error. In three error
cases the generated prose also contradicts the structured label, identifying a
separate generation-consistency issue.

The direct-claim counterexample subset yields 95\% accuracy for Full and 85\%
for Causal Path RAG. However, all 20 samples are evaluated by direct-edge
topology and none executes $do(X_M=0)$. This subset supports the value of
query-aware direct-claim checking, not the effectiveness of structural mediator
intervention.

\subsection{Retrieval, Path Ranking, and Runtime}

Causal Path RAG and Full share the same retrieval output. Article Recall@5,
Rule Recall@5, and Event Recall@5 are 75.40\%, 71.20\%, and 74.79\%.
Top-1 Exact Path is 38.56\%, while Oracle Exact Path is 75.42\%. The 36.86-point
gap shows that the gold path is often in the candidate set but ranked below the
first position; retrieval/path equality also localizes the Full--Causal
decision difference to the selective verification layer rather than retrieval.

Full requires 3.324 seconds per successful question versus 2.107 seconds for
Causal Path RAG, a $1.58\times$ average-time increase. The measured trade-off is
therefore a statistically reliable 6.0-point decision gain at approximately
58\% additional per-sample runtime.

% =========================================================
% VII. DISCUSSION AND LIMITATIONS
% =========================================================
\section{Discussion and Limitations}
\label{sec:discussion_limitations}

The controlled experiment shows that structured event--rule retrieval and
selective verification improve exact decision labels. Full LegalSCM reaches
96.80\% accuracy versus 90.80\% for Causal Path RAG, with a statistically
significant paired gain. Since both systems expose identical retrieval and path
outputs, and all 17 Full-only gains occur after verifier commitment, the
difference is localized to committed factual validation and grounded
direct-edge checks.

The result does not imply a uniform quality gain. Answer F1 and ROUGE-L are
slightly lower for Full, and their paired differences, as well as the Citation
F1 difference, are not significant. Structured-label correctness and lexical
answer overlap therefore remain distinct objectives. Three erroneous cases
also contain prose that conflicts with the structured decision, motivating an
explicit generation-consistency constraint.

The closed-book LLM result is strongly prompt- and task-dependent. Its 0\%
accuracy arises because all 250 successful outputs abstain as
\textsc{uncertain}, while the gold set contains no abstention class. It is not
evidence that Qwen3-8B has zero legal knowledge. The 230/20 gold-label imbalance
also makes raw accuracy easy to inflate, so Balanced Accuracy, Macro-F1, class
coverage, and the confusion matrix should accompany it.

Most importantly, the reported experiment does not validate structural
mediator intervention. The 20 negative samples are direct-edge counterexamples,
and the saved Full run executes zero $do(X_M=0)$ operations. The 95\% subset
score measures grounded direct-edge topology. The hard-intervention mechanism
is implemented and formally separated from $G^{-M}$ node deletion, but requires
a new benchmark with explicit removal/necessity queries before effectiveness
claims can be made. Likewise, the engine supports conjunctive bodies and
alternative disjunctive mechanisms, whereas the current imported rules are
unary positive rules; rich AND/OR behavior is not empirically tested.

Top-1 Exact Path remains 38.56\% versus 75.42\% Oracle Exact Path. Seven of the
eight Full errors occur after a missing-primary-path fallback, confirming that
first-hop retrieval and path ranking remain bottlenecks. In addition, the
benchmark is small, graph-derived, and built from the same structured resource
used by the system. Manual expert review, out-of-graph legal questions, and
external-domain testing are needed for external validity.

Finally, LegalSCM depends on rule extraction, event normalization, endpoint and
mediator grounding, and graph construction. It is a deterministic normative
rule model, not causal discovery or effect estimation. Its $do(\cdot)$ operator
is meaningful only relative to encoded mechanisms and lacks the exogenous
abduction layer required for full unit-level counterfactual analysis.

% =========================================================
% VIII. CONCLUSION
% =========================================================
\section{Conclusion}
\label{sec:conclusion}

LegalSCM-RAG combines causal-path retrieval, three-valued rule inference, and a
selective commit gate for Vietnamese multi-hop legal QA. On the current
250-question benchmark, it improves decision accuracy over Causal Path RAG by
6.0 points with paired statistical support, but does not significantly improve
answer or citation metrics.

The strongest supported conclusion is therefore that selective factual and
direct-claim verification improves decision reliability in this controlled
setting. Future work will construct explicit mediator-intervention and
multi-antecedent rule benchmarks, improve first-hop retrieval and endpoint
grounding, enforce answer--decision consistency, and obtain expert evaluation
on legal questions not generated from the source graph.


\bibliographystyle{splncs04}
\bibliography{../references}

\end{document}
